% Chapter 1: Introduction

\section{Problem and Contribution}

Scientific fields are usually described by names, topics, citation links, publication counts, or positions on maps. These views are useful, but they do not directly measure how the papers of a field are arranged. Two fields can belong to different OpenAlex domains and still have similarly shaped paper distributions; two fields can be close in semantic location and differ sharply in dispersion, local density, hubness, or spectral structure.

This thesis treats scientific fields as distributions of embedded documents. I use OpenAlex title--abstract records, represent them with SPECTER2 document embeddings, and measure field morphology in the original 768-dimensional embedding space. By embedding-space morphology I mean the measurable structural shape of a field's paper distribution: its spread around a center, its local neighborhoods, its hub concentration, and its principal directions of variation.

The active morphology profile has eight structural metrics. Temporal evolution is measured by recomputing those same eight metrics in five-year windows. Net semantic displacement is handled separately through early--late centroid drift: it describes how far a subfield's semantic center moves, not the structural shape used in static profiles, PCA, clustering, or morphological distances.

The contribution is a reproducible framework for separating five related but distinct objects of analysis: labels, semantic locations, static shapes, temporal trajectories, and relations between fields. The thesis does not argue that morphology replaces taxonomies, citation maps, topic models, or visual maps. It argues that those tools leave a measurable layer underdescribed: the internal arrangement of papers once they have shared document coordinates.

\section{Research Questions}

The empirical analysis is organized around three research questions:

\begin{enumerate}
    \item How can the internal structure of scientific fields be measured using scientific document embeddings?
    \item How does the internal structure of scientific fields vary across disciplines and evolve over time?
    \item Which scientific fields converge or diverge in their structural profiles over time?
\end{enumerate}

These questions are descriptive rather than causal. A dispersed field is not better than a compact one, a dynamic field is not necessarily more important, and a low morphological distance does not mean that two fields study the same topic. The goal is to define and test a measurement language for field shape.

\section{Roadmap}

Chapter 2 positions the thesis in science mapping and scientific document representation. Chapter 3 describes the corpus, SPECTER2 representation, structural morphology metrics, temporal windows, and workflow. Chapter 4 compares static morphology across domains, fields, and subfields; Chapter 5 examines temporal evolution, dynamic profiles, and centroid drift as a separate displacement indicator; Chapter 6 studies morphological similarity, taxonomy, convergence, and divergence. Chapters 7 and 8 discuss the contribution, limits, and final findings.
